\section{模型检验与误差分析}\label{sec_validation}
\subsection{水分收支与有效热方程检验}
固定域对半离散方程求和，内部通量相消，得到
\begin{equation}\label{eq_mass_fixed}
\sum_iW_i[C_i(t)-C_i(0)]
=-Rh_m\int_0^t[C_s(\tau)-C_a(\tau)]\dd\tau.
\end{equation}
式\eqref{eq_mass_fixed}按参考干密度归一化，不能直接以kg标记；收缩域使用式\eqref{eq_mass_shrink}。

问题一热容量恒定，总有效蓄热增量应与累计对流输入一致。问题二、三则令
\[
E(t)=2\pi L\sum_iW_iB(C_i)(T_i-T_{\mathrm{ini}}),
\]
由式\eqref{eq_B_identity}推导
\begin{equation}\label{eq_heat_check}
\begin{gathered}
E(t)-E(0)-2\pi L\int_0^t\sum_iW_iB'(C_i)(T_i-T_{\mathrm{ini}})\dot C_i\dd\tau\\
=2\pi LRh\int_0^t(T_a-T_s)\dd\tau.
\end{gathered}
\end{equation}
组成修正项不能省略。相对残差用收支差除以累计交换量尺度，近零时设正下限；第一问定义见式\eqref{eq_q1_balance_residual}。第三问在BDF接受步连续解上作4点、8点Gauss积分对照。表\ref{tab_balance}汇总结果，第四问热量仅检验加权半离散方程及温度收敛。
\begin{table}[htbp]\centering\small
\caption{各问收支检验；相对残差仅针对所选有效模型}\label{tab_balance}
\begin{tabular}{cccc}\toprule
问题 & 水分相对残差 & 有效热收支相对残差 & 检验范围\\\midrule
一 & $2.4254\times10^{-11}$ & $2.0484\times10^{-9}$ & 0至1800 s\\
二 & $4.312\times10^{-7}$ & $3.565\times10^{-8}$ & 0至10800 s\\
三 & $4.0221\times10^{-11}$ & $3.6640\times10^{-8}$ & 至报告时刻\\
四 & $4.0558\times10^{-9}$ & 不作同类累计比较 & 基准计算至72 h\\\bottomrule
\end{tabular}
\end{table}
各问残差均远低于$10^{-6}$的检查阈值，说明离散通量与所建方程一致；该检验针对本文的有效模型，不涉及潜热、相平衡或端面简化所对应的物理误差。

\subsection{空间与时间加密}
在相同物理时刻比较固定域共享节点及收缩域共享材料节点，以各配置的节点最大绝对差为指标。表\ref{tab_convergence}分别列出温度和含水率结果。
\begin{table}[htbp]\centering\small
\caption{主要网格与时间加密差异}\label{tab_convergence}
\begin{tabular}{clrrrr}\toprule
问题 & 空间比较 & $\Delta T_{\rm space}$ & $\Delta C_{\rm space}$ & $\Delta T_{\rm time}$ & $\Delta C_{\rm time}$\\
 & & $^\circ$C & kg/kg & $^\circ$C & kg/kg\\\midrule
一 & $G_4/G_8$ & $5.895\!\times\!10^{-7}$ & $4.564\!\times\!10^{-6}$ & $2.874\!\times\!10^{-8}$ & $1.179\!\times\!10^{-9}$\\
二 & $G_2/G_4$ & $2.452\!\times\!10^{-6}$ & $1.493\!\times\!10^{-5}$ & $4.108\!\times\!10^{-7}$ & $2.381\!\times\!10^{-8}$\\
三 & $G_8/G_{16}$ & $9.545\!\times\!10^{-7}$ & $1.537\!\times\!10^{-5}$ & $1.002\!\times\!10^{-6}$ & $7.716\!\times\!10^{-9}$\\
四 & $G_8/G_{16}$ & $1.012\!\times\!10^{-5}$ & $8.681\!\times\!10^{-7}$ & $1.188\!\times\!10^{-5}$ & $6.383\!\times\!10^{-8}$\\\bottomrule
\end{tabular}
\end{table}
前两问覆盖全部规定输出点，后两问还覆盖内部节点，第四问比较至72 h。第二问另以$G_4/G_8$复核，温度、含水率最大差为$6.244\times10^{-7}$ $^\circ$C、$3.732\times10^{-6}$ kg/kg，复现方式见支撑材料。六张题定表在相应加密下四位小数一致；表中差值为不同计算配置之间的比较结果，不构成真解的误差上界。各配置下的临界时间见表\ref{tab_event_convergence}。
\begin{table}[htbp]\centering\small
\caption{不同计算配置下的临界烘干时间（h）}\label{tab_event_convergence}
\begin{tabular}{crr}\toprule
配置 & 问题三 & 问题四\\\midrule
$G_4$ & 57.4731739851 & 51.0872487284\\
$G_8$ & 57.4725195022 & 51.0871229810\\
$G_{16}$ & 57.4723569654 & 51.0870914570\\
$G_8$时间收紧 & 57.4725194974 & 51.0871221801\\\bottomrule
\end{tabular}
\end{table}

\subsection{独立对照与输出一致性}
问题一的Bessel--Duhamel温度对照及冻结扩散系数辅助测试见第\ref{sec_q1_validation}节。问题二在短时制造解测试中保留变系数导数项，温度和含水率最大误差分别为$6.0949\times10^{-8}$ $^\circ$C和$8.9808\times10^{-9}$ kg/kg。制造解用于检验实现，不是药材实验数据。

第三问与第二问在1800、3600、7200、9000 s的重叠输出比较中，温度、水分最大差分别为$4.66\times10^{-7}$ $^\circ$C和$9.12\times10^{-7}$ kg/kg，说明独立起算的两个过程在共享区间内一致。第四问则采用同一附录4物性的固定半径退化检查及关闭边界交换的纯收缩检查，以保证退化参照与本问物性一致。

论文表格、Excel结果与源数组均由同一未舍入数值解生成，并另行核对第三问的准确终时、第四问的当前半径及域外标记。上述检查针对数值实现，不替代物理验证。

\subsection{同物性固定半径对照的检验}
固定组通过独立物理环层算子检查。空间加密和时间收紧的临界时间差为0.498231 s和0.013113 s；至132 h共同保存时刻的温度、水分空间差分别为$1.3522\times10^{-5}$ $^\circ$C、$1.9052\times10^{-6}$ kg/kg，时间差分别为$1.4730\times10^{-5}$ $^\circ$C、$6.4161\times10^{-8}$ kg/kg。三组累计水分相对残差均低于$10^{-6}$。129.8448 h时三组最大含水率均低于0.15，生产值加经验裕量$1.9437\times10^{-7}$后仍达标，支持表\ref{tab_q4_control}的数值可比性。

\subsection{局部敏感性与适用边界}
式\eqref{eq_logD}给出扩散系数对状态变量的局部敏感性。在$T=50\ ^\circ$C、$C=0.15$时，问题三有
\[
\frac{\partial\ln D}{\partial T}\approx0.0369\ \mathrm{K^{-1}},\qquad
\frac{\partial\ln D}{\partial C}=20\ \mathrm{(kg/kg)^{-1}}.
\]
局部升温1 $^\circ$C对应$D$约3.69\%的一阶变化，不能等同于时长变化。问题四有$\partial\ln D/\partial U=0.30/U^2$，阈值处为13.3333 $(\mathrm{kg/kg})^{-1}$；几何因子满足$\partial\ln(R^{-2})/\partial\ln R=-2$，其局部方向仍需全过程对照检验。

若事件附近最大值位置不变且$M'(t_*)\ne0$，一阶扰动给出
\begin{equation}\label{eq_event_sensitivity}
\delta t_*\approx-\frac{\delta M(t_*)}{M'(t_*)}.
\end{equation}
后期曲线越平缓，相同含水率误差越可能放大为时间偏差，故需直接检验事件收敛。表\ref{tab_q3_sensitivity}的有限输入扰动量化了指定情景的时长变化，不构成全局稳健性保证。

第三问径向近似的二维端面校核已在第\ref{sec_axisymmetric}节给出；其作用是检验本题条件下的维度简化，与本章的离散收支和数值加密检验相互补充。
\FloatBarrier
